\section{Threat Model and Safety Considerations}

This section outlines the assumptions regarding the adversary, the scope of targeted threats, and the safety protocols implemented to ensure secure experimentation.

\subsection{Adversarial Assumptions}

We focus on ransomware that has successfully established a presence on a host system. Specifically, the adversary is assumed to possess the capability to execute malicious code following an initial compromise, invoke standard system APIs and cryptographic libraries for file encryption, and perform high-volume file operations such as reads, writes, and renames with anomalous frequency. Additionally, the attacker may spawn new processes, attempt privilege escalation, establish persistence mechanisms, and optionally communicate with a command-and-control server for key exchange or status reporting.

The adversary model encompasses several ransomware categories based on their operational characteristics. \emph{Crypto-ransomware} focuses primarily on file encryption, employing hybrid cryptographic schemes combining symmetric and asymmetric algorithms to ensure data remains inaccessible without payment. \emph{Locker ransomware} restricts system access without necessarily encrypting files, typically targeting system boot processes or user interface components. \emph{Wiper variants} masquerade as ransomware but prioritize data destruction over extortion, often employed in state-sponsored or destructive attacks. ADRDS is primarily designed to detect crypto-ransomware due to its prevalence and the distinctive behavioral signatures associated with encryption operations.

\subsection{Attack Lifecycle Phases}

The ransomware attack lifecycle comprises distinct phases, each exhibiting characteristic behavioral patterns. Understanding these phases informs the detection strategy and enables stage-aware alerting.

\textbf{Initial Access and Execution.} Following successful infection through vectors such as phishing, exploit kits, or supply chain compromise, the ransomware payload executes on the target system. This phase may involve unpacking, decryption of embedded components, and initial environment reconnaissance.

\textbf{Privilege Escalation.} Attackers frequently attempt to elevate privileges to gain broader system access. This may involve exploiting local vulnerabilities, abusing misconfigured services, or leveraging credential theft techniques. Elevated privileges enable encryption of protected system files and facilitate persistence.

\textbf{Persistence Establishment.} To survive system reboots and ensure continued access, ransomware often establishes persistence through registry modifications, scheduled tasks, or service installation. These activities generate observable artifacts that contribute to detection.

\textbf{Lateral Movement.} In enterprise environments, ransomware may propagate to additional systems through network shares, remote execution protocols, or exploitation of networked services. This phase amplifies the attack impact and complicates remediation.

\textbf{Data Exfiltration.} Modern ransomware increasingly incorporates data theft prior to encryption, enabling double extortion tactics. This phase involves file identification, compression, and transmission to attacker-controlled infrastructure.

\textbf{Encryption and Ransom Demand.} The final operational phase involves the systematic encryption of identified files and the presentation of ransom demands. This phase exhibits the most distinctive behavioral signatures, including high-volume file operations and entropy anomalies.

\textbf{Defender Visibility.}
From a defensive perspective, it is assumed that the monitoring system maintains visibility over host-level telemetry, including file system events, process activity, API calls, and resource utilization. ADRDS specifically targets the pre-encryption and active encryption phases, as these periods exhibit behavioral divergence from standard operational workloads.

\textbf{Scope Limitations.}
Certain vectors are explicitly excluded from the scope of this study. Initial infection vectors, such as phishing campaigns, exploit kits, or supply-chain compromises, are outside the detection model. Furthermore, post-encryption activities, including ransom negotiation, payment infrastructure, and monetization strategies, are not addressed.

\subsection{Safety Protocols}

\textbf{Simulation-Based Validation.}
ADRDS avoids the execution of live ransomware binaries. Instead, it utilizes parameterized behavioral simulations to replicate indicators of compromise—such as entropy spikes, rapid file writes, and anomalous process activity—without executing actual encryption or destructive payloads. This design ensures that demonstrations and experiments remain inherently safe.

The simulation engine implements precise control over behavioral parameters, enabling the reproduction of specific ransomware families and attack scenarios. Parameters include encryption velocity, file targeting patterns, process spawning behavior, and resource utilization profiles. This granularity supports both research requiring specific attack characteristics and educational demonstrations requiring progressive threat escalation.

\textbf{Environment Isolation.}
All experiments are conducted within sandboxed, air-gapped environments, devoid of connectivity to production systems or sensitive networks. This isolation mitigates the risk of accidental data exposure and prevents the lateral movement of simulated threats. Network interfaces are disabled or restricted to local communication, and file system access is confined to designated experimental directories.

\textbf{Ethical Considerations.}
ADRDS is developed strictly for defensive research, education, and demonstration purposes. The system does not contain exploit code or instructional material susceptible to weaponization. Users are expected to adhere to applicable legal frameworks, institutional policies, and ethical guidelines when utilizing the platform. The development team maintains responsible disclosure practices and engages with the security research community to ensure appropriate use.

\subsection{Privacy and Data Handling}

The system does not collect, store, or transmit personally identifiable information (PII). All telemetry utilized for detection and visualization is either synthetically generated or derived from non-sensitive experimental datasets, ensuring alignment with privacy-preserving research standards. Data retention policies limit storage duration, and all experimental data remains local to the execution environment.
